\section*{Day 4}
\section*{Exercise 5}
In this exercise we explore variance reduction methods. Parts 1--4 estimate the integral $\int_0^1 e^x\,dx$ (whose exact value is $e-1 \approx 1.718282$) using four different techniques, all with comparable computer resources, aiming for 100 function evaluations for each method. Part 5 applies these ideas to the blocking system from Exercise 4, and lastly parts 7-9 deal with importance sampling.

Throughout, we report the $95\%$ confidence interval (CI) width as our measure of estimator precision. This is equivalent to reporting variance: the CI width is $2\,t_{0.975}(n-1)\,s/\sqrt{n}$, so it is a direct function of the sample variance $s^2$ and the number of observations $n$. Since both $s^2$ and $n$ differ across methods, we believe the CI width is the single most interpretable summary — it already accounts for both sources of precision.

\subsection*{Part 1 -- Crude Monte Carlo}
We use the crude Monte Carlo estimator to estimate the integral
\[
\int_0^1 e^x \, dx
\]
We draw $n$ samples $U_i \sim \operatorname{Uniform}(0,1)$ and estimate the integral as the sample mean of $e^{U_i}$,
\[
\frac{1}{n}\sum_{i=1}^n e^{U_i},
\]
with a $95\%$ confidence interval obtained from the $t$-distribution. With $n=100$ we get
\begin{center}
mean $= 1.714466$, \quad $95\%$ CI $= (1.617300,\ 1.811632)$.
\end{center}
The true value $e-1 = 1.718282$ lies inside the interval, as expected.

\subsection*{Part 2 -- Antithetic variables}
To reduce the variance we use antithetic sampling. For each $U$ we also use its antithetic counterpart $1-U$ and average the two:
\[
\frac{e^{U}+e^{1-U}}{2}.
\]
Since $U$ and $1-U$ have the same distribution, $e^{U}$ and $e^{1-U}$ also do, so $\mathbb{E}\!\left[\tfrac{1}{2}(e^{U}+e^{1-U})\right] = e-1$ and the estimator remains unbiased. Because $e^{U}$ and $e^{1-U}$ are negatively correlated, their average has reduced variance. Using $n=50$ pairs (i.e. $100$ function evaluations, comparable to Part 1) we get
\begin{center}
mean $= 1.711206$, \quad $95\%$ CI $= (1.692801,\ 1.729611)$.
\end{center}

Notice that the confidence interval is significantly tighter now, even with the sample size going from $n=100$ in part 1 to $n=50$ here. The confidence interval went from CI = (1.617300, 1.811632) which has a width of  0,1943 in part 1 to CI = (1.692801, 1.729611) which has width = 0,0368 in part 2 indicating that variance reduction was successfully achieved by incorporating the antithetic variable.

\subsection*{Part 3 -- Control variate}
Next we use a control variate. This is a technique where we add a related random variable, with its mean subtracted so that the added term $Z-\mathbb{E}[Z]$ satisfies $\mathbb{E}\!\left[Z-\mathbb{E}[Z]\right]=0$ and the estimator thus remains unbiased. The idea is then to scale this term by a constant $c$ chosen to maximize the variance reduction (i.e.\ minimize the overall variance). The final added term is thus $c\,(Z-\mathbb{E}[Z])$. The optimal value for the constant is given as
\[
c = -\frac{\operatorname{Cov}(X_i, Z_i)}{\operatorname{Var}(Z_i)},
\]
where $X$ is the target and $Z$ is the control variate. We choose $Z = U$, since it is strongly correlated with $e^{U}$ and is very simple to evaluate. In the simulation $c$ is estimated empirically from the sample by
\begin{lstlisting}[language=Python, basicstyle=\ttfamily\small]
c = -np.cov(X, Z, ddof=1)[0, 1] / np.var(Z, ddof=1)
\end{lstlisting}
which gives $c \approx -1.688$.

We use $n=100$, since there is only one function evaluation per iteration, keeping
the computational resources comparable to the previous parts. Using this control
variate approach we obtain
\begin{center}
mean $= 1.716209$, \quad $95\%$ CI $= (1.704551,\ 1.727866)$.
\end{center}
The width of the CI is thus further reduced compared to Part 2 (antithetic variables), and the interval still contains the true value, as expected. It is worth noting that both methods have essentially the same theoretical per-observation variance ($\approx 0.0039$), so the improvement here is partly due to $n$: control variates yield $100$ independent observations from $100$ evaluations, whereas antithetic variables yield only $50$ pairs. The methods are not quite on equal footing in that sense, even though the total number of function evaluations is the same.

\subsection*{Part 4 -- Stratified sampling}
Stratified sampling is a simple idea: instead of drawing $U$ uniformly on $[0,1)$, we split $[0,1)$ into $k$ subintervals of equal length and draw one sample from each, repeating $n/k$ times so that we still obtain $n$ total samples. 
Within stratum $j$ the draw is confined to that subinterval,
\[
U_j = \frac{j+V}{k}, \qquad V \sim \operatorname{Uniform}(0,1).
\]

This reduces variance because the samples can no longer cluster by chance in one region of the domain. There is now no variance between different strata but only some possible variance within the strata itself.
Now each iteration of the loop results in one estimator value:
\[
Y_i = \frac{1}{k}\sum_{j=1}^{k} e^{U_j},
\]
We repeat this $n/k$ times and thus get $n/k$ estimators. Therefore the mean and confidence interval are calculated based on this composite estimator. 
With $n=100$ and $k=10$ we get
\begin{center}
mean $= 1.722733$, \quad $95\%$ CI $= (1.715688,\ 1.729779)$.
\end{center}

Once again, we managed to reduce the width of the confidence interval even below the confidence interval from part 2 and 3, indicating that stratified sampling works well for variance reduction, at least for this problem.

\subsection*{Comparison of Parts 1--4.}
Table~\ref{tab:vr14} compares the confidence interval widths of the four estimators at comparable computer resources. All three variance reduction methods clearly outperform crude Monte Carlo, and out of our methods and approaches, stratified sampling gives the narrowest interval, followed by the control variate and then antithetic sampling.

\begin{table}[h]
\centering
\caption{$95\%$ CI widths for estimates of $\int_0^1 e^x\,dx$ at comparable compute.}
\label{tab:vr14}
\begin{tabular}{lrr}
\toprule
Method & CI width & Reduction vs.\ crude \\
\midrule
Crude (Part 1)        & $0.1943$ & --- \\
Antithetic (Part 2)   & $0.0368$ & $\sim 5\times$ \\
Control variate (Part 3) & $0.0233$ & $\sim 8\times$ \\
Stratified (Part 4)   & $0.0141$ & $\sim 14\times$ \\
\bottomrule
\end{tabular}
\end{table}

\begin{figure}
    \centering
    \includegraphics[width=0.7\linewidth]{plots//day4/ci_comparison.png}
    \caption{95\% confidence intervals for $\int_0^1 e^x\,dx$ across the four methods at comparable compute (100 function evaluation each). The red dashed line marks the true value $e-1$.}
    \label{fig:ci_comparison}
\end{figure}

\subsection*{Part 5 -- Control variates on the blocking system}
We now apply variance reduction to the blocking system from Exercise 4, using a control variate to reduce the variance of the estimated blocking fraction.

We use the sum of the inter-arrival times $S$ as the control variate. Its expectation is known: with $n$ customers and mean inter-arrival time $1$ we have $\mathbb{E}[S] = n$. Intuitively, when arrivals are spread out (large $S$) the system is less loaded and fewer customers are blocked, so $S$ is negatively correlated with the blocking fraction $B$ and is a natural control variate.

We first perform $100$ runs to estimate the optimal coefficient
\[
c = \frac{\operatorname{Cov}(B, S)}{\operatorname{Var}(S)},
\]
and then run a fresh batch of simulations to find the estimate $Y_i = B_i - c\,(S_i - n)$:

\begin{lstlisting}[language=Python, basicstyle=\ttfamily\small]
Bs, Ss, Ys = [], [], []
n = 10000
n_reps = 100
for i in range(n_reps):
    arrival_fn = lambda: np.random.exponential(1.0)
    service_fn = lambda: np.random.exponential(8.0)
    B_i, S_i = simulate_cv(arrival_fn, service_fn, m=10, n=n)
    Bs.append(B_i)
    Ss.append(S_i)
    Ys.append(B_i - c * (S_i - n))

mean = np.mean(Ys)
std = np.std(Ys, ddof=1)
half_width = scipy.stats.t.ppf(0.975, df=n_reps-1) * std / np.sqrt(n_reps)
\end{lstlisting}

The results of the simulation run are shown in Table~\ref{tab:part5_results}, with the Erlang B reference value $0.121661$ for comparison. The control variate reduces the CI width by about $20\%$. The reduction is less than what we saw in parts 1-4, but it is still a reasonable reduction in variance for no extra per-simulation cost (beyond a short pilot run to estimate the $c$ constant).

\begin{table}[h]
\centering
\caption{Means and CI widths for the simulated blocking system.}
\label{tab:part5_results}
\begin{tabular}{lrrr}
\toprule
Method & Mean & CI width & Reduction vs.\ crude \\
\midrule
Crude (no CV)   & $0.121796$ & $0.002216$ & --- \\
Control variate & $0.121185$ & $0.001767$ & $\sim 1.25\times$ \\
\bottomrule
\end{tabular}
\end{table}

\subsection*{Part 7 -- Importance sampling for $P(Z > a)$} We estimate $P(Z > a)$ for $Z \sim N(0,1)$. First we use crude Monte Carlo, and then importance sampling with a proposal density that is normal with mean $a$ and variance $\sigma^2$, for different values of $a$, $\sigma^2$, and sample size $n$. The results are shown in Table~\ref{tab:is7}.

\begin{table}[h]
\centering
\caption{Crude vs.\ importance sampling for $P(Z>a)$. True values:
$P(Z>2) = 2.28\times10^{-2}$, $P(Z>4) = 3.17\times10^{-5}$.}
\label{tab:is7}
\begin{tabular}{cclrr}
\toprule
$n$ & $a$ & Method & Mean & CI width \\
\midrule
\multirow{8}{*}{$100$}
 & \multirow{4}{*}{$2$} & Crude            & $3.00\times10^{-2}$ & $6.80\times10^{-2}$ \\
 &                      & IS, $\sigma^2=1$ & $2.15\times10^{-2}$ & $1.37\times10^{-2}$ \\
 &                      & IS, $\sigma^2=2$ & $2.10\times10^{-2}$ & $1.47\times10^{-2}$ \\
 &                      & IS, $\sigma^2=3$ & $1.56\times10^{-2}$ & $1.36\times10^{-2}$ \\
\cmidrule(l){2-5}
 & \multirow{4}{*}{$4$} & Crude            & $0.00$              & $0.00$ \\
 &                      & IS, $\sigma^2=1$ & $3.67\times10^{-5}$ & $3.16\times10^{-5}$ \\
 &                      & IS, $\sigma^2=2$ & $2.40\times10^{-5}$ & $2.64\times10^{-5}$ \\
 &                      & IS, $\sigma^2=3$ & $3.09\times10^{-5}$ & $3.12\times10^{-5}$ \\
\midrule
\multirow{8}{*}{$1000$}
 & \multirow{4}{*}{$2$} & Crude            & $2.10\times10^{-2}$ & $1.78\times10^{-2}$ \\
 &                      & IS, $\sigma^2=1$ & $2.41\times10^{-2}$ & $4.36\times10^{-3}$ \\
 &                      & IS, $\sigma^2=2$ & $2.10\times10^{-2}$ & $5.03\times10^{-3}$ \\
 &                      & IS, $\sigma^2=3$ & $2.32\times10^{-2}$ & $6.24\times10^{-3}$ \\
\cmidrule(l){2-5}
 & \multirow{4}{*}{$4$} & Crude            & $0.00$              & $0.00$ \\
 &                      & IS, $\sigma^2=1$ & $3.03\times10^{-5}$ & $7.92\times10^{-6}$ \\
 &                      & IS, $\sigma^2=2$ & $2.96\times10^{-5}$ & $1.02\times10^{-5}$ \\
 &                      & IS, $\sigma^2=3$ & $3.16\times10^{-5}$ & $1.11\times10^{-5}$ \\
\midrule
\multirow{8}{*}{$10000$}
 & \multirow{4}{*}{$2$} & Crude            & $2.22\times10^{-2}$ & $5.78\times10^{-3}$ \\
 &                      & IS, $\sigma^2=1$ & $2.22\times10^{-2}$ & $1.35\times10^{-3}$ \\
 &                      & IS, $\sigma^2=2$ & $2.33\times10^{-2}$ & $1.72\times10^{-3}$ \\
 &                      & IS, $\sigma^2=3$ & $2.26\times10^{-2}$ & $1.91\times10^{-3}$ \\
\cmidrule(l){2-5}
 & \multirow{4}{*}{$4$} & Crude            & $0.00$              & $0.00$ \\
 &                      & IS, $\sigma^2=1$ & $3.10\times10^{-5}$ & $2.60\times10^{-6}$ \\
 &                      & IS, $\sigma^2=2$ & $3.17\times10^{-5}$ & $3.21\times10^{-6}$ \\
 &                      & IS, $\sigma^2=3$ & $3.19\times10^{-5}$ & $3.63\times10^{-6}$ \\
\bottomrule
\end{tabular}
\end{table}

It is remarkable how much better importance sampling performs. Already at $n=100$ it captures the rare events and produces a much narrower confidence interval than crude Monte Carlo (e.g. $1.47\times10^{-2}$ vs.\ $6.80\times10^{-2}$ for $a=2$). So for essentially the same compute, importance sampling gives far lower variance.

The effect is even more dramatic for the rare event $a=4$: crude Monte Carlo returns exactly $0$ for $n \le 1000$ because it never draws a sample beyond $4$, whereas importance sampling estimates the true probability $3.17\times10^{-5}$ relatively accurately already at $n=100$ and very accurately at the larger sample sizes. Finally, if performance is measured by confidence interval width, the smallest variance $\sigma^2 = 1$ performed best overall among the proposals tried, having the lowest CI width in 5 out of 6 comparisons.

\subsection*{Part 8 -- Importance sampling with an exponential proposal}
We estimate the integral
\[
\theta = \int_0^1 e^x \, dx
\]
using importance sampling. Following the notation from the lecture,
we write the quantity of interest as
\[
\theta = \mathbb{E}[h(X)] = \int_{-\infty}^{\infty} h(x)\,f(x)\,dx,
\qquad
h(x) = e^x,
\qquad
f(x) = \mathbf{1}_{\{0 \le x \le 1\}},
\]
where $h$ is the integrand and $f$ is the $\operatorname{Uniform}(0,1)$ density. The indicator makes
$f$ zero outside $[0,1]$, so this reproduces the original integral.

Instead of sampling from $f$, we sample from the proposal density
\[
g(x) = \lambda e^{-\lambda x}\]
Since $g(x) > 0$ wherever $f(x) > 0$, the identity from slide~27 gives
\[
\theta = \int h(x)\frac{f(x)}{g(x)}\,g(x)\,dx
       = \mathbb{E}\!\left[h(X)\frac{f(X)}{g(X)}\right],
\]
with estimator $\hat{\theta}_{IS} = \tfrac{1}{n}\sum_{i=1}^n h(X_i)\frac{f(X_i)}{g(X_i)}$.
The importance sampling random variable is therefore
\[
W = h(X)\frac{f(X)}{g(X)}
  = \frac{e^X\,\mathbf{1}_{\{0\le X\le 1\}}}{\lambda e^{-\lambda X}}
  = \frac{1}{\lambda}e^{(1+\lambda)X}\mathbf{1}_{\{0\le X\le 1\}}.
\]
Since $X \sim \operatorname{Exp}(\lambda)$ is always non-negative,
$\mathbf{1}_{\{0\le X\le 1\}} = \mathbf{1}_{\{X\le 1\}}$.

To find the optimal $\lambda$ we analyze $\operatorname{Var}(W)$. Since the estimator is unbiased, $\mathbb{E}[W]=\int_0^1 e^x\,dx = e-1$, so
\[
\operatorname{Var}(W)
=
\mathbb{E}[W^2]-\mathbb{E}[W]^2
=
\mathbb{E}[W^2]-(e-1)^2 .
\]
Now we calculate $\mathbb{E}[W^2]$
\[
\mathbb{E}[W^2]
=\int_{-\infty}^\infty W(x)^2 g(x)\,dx
=\int_0^1 \frac{e^{2x(1+\lambda)}}{\lambda^2}\,\lambda e^{-\lambda x}\,dx
=\frac{1}{\lambda}\int_0^1 e^{(2+\lambda)x}\,dx ,
\]
hence
\[
\mathbb{E}[W^2]
=\frac{1}{\lambda}\left[\frac{e^{(2+\lambda)x}}{2+\lambda}\right]_0^1
=\frac{e^{2+\lambda}-1}{\lambda(2+\lambda)},
\qquad
\operatorname{Var}(W)
=\frac{e^{2+\lambda}-1}{\lambda(2+\lambda)}-(e-1)^2 .
\]
The variance of the estimator is $\operatorname{Var}(\hat{I}_{IS}) =
\tfrac{1}{n}\operatorname{Var}(W)$. Since $1/n$ and $(e-1)^2$ is independent of $\lambda$,
minimising the variance is equivalent to minimising
$\frac{e^{2+\lambda}-1}{\lambda(2+\lambda)}$. Done numerically, this gives
\footnote{ChatGPT-5.5 was used to assist with the derivation.}

\[
\boxed{\lambda^* \approx 1.3548}
\]

with minimum variance $\operatorname{Var}(W) \approx 3.1288$.

Running the simulation gives the results in Table~\ref{tab:is8}.

\begin{table}[h]
\centering
\caption{Crude MC vs.\ exponential importance sampling for $\int_0^1 e^x\,dx$
($n=10{,}000$). True value $e-1 = 1.718282$.}
\label{tab:is8}
\begin{tabular}{lrr}
\toprule
Estimator & Mean & CI width \\
\midrule
Crude MC                       & $1.717034$ & $0.019294$ \\
IS, $\lambda=0.500$            & $1.645545$ & $0.093913$ \\
IS, $\lambda=1.000$            & $1.693668$ & $0.072518$ \\
IS, $\lambda^*=1.355$          & $1.732042$ & $0.070012$ \\
IS, $\lambda=2.000$            & $1.710978$ & $0.075783$ \\
IS, $\lambda=5.000$            & $1.629328$ & $0.196602$ \\
\bottomrule
\end{tabular}
\end{table}

The simulation confirms the analysis: among the exponential proposals, the narrowest interval is indeed obtained at $\lambda^* \approx 1.355$, with CI width $0.070$. However, this is still about $3.6\times$ \emph{wider} than crude Monte Carlo ($0.019$). This matches the analytical result, $\operatorname{Var}(W) \approx 3.13$ at $\lambda^*$ versus a crude MC variance of only $\approx 0.242$. So importance sampling with an exponential distribution does \emph{not} reduce the variance here. The reason is that the exponential density is decreasing on $[0,1]$, whereas the integrand $e^x$ is increasing, so the proposal is the wrong shape and over-weights the region where the integrand is smallest.

\subsection*{Part 9 -- Pareto first-moment distribution}